% R008 — Guide to Cultivating Complex Analysis (Bahasa Indonesia)
% Reader-facing translation unit: source/ca-v1.9/ca.tex lines 761–786.
% Source release: v1.9 (commit a4d25d9ba37a486a5a020219eb8bd4e6602fd14c).
% Translation provenance: OpenAI Codex gpt-5.6-sol, Ultra, acting on the user's request.
% Preserve source labels, glossary hooks, mathematical notation, and structure.

Ketika kita menulis bilangan real $x$, kita mengidentifikasikannya dengan
bilangan kompleks $(x,0)$.
Dengan identifikasi ini, $\R \subset \C$.
Kita juga mendefinisikan \emph{\myindex{satuan imajiner}}%
\footnote{Waspadalah terhadap para insinyur; mereka mengira satuan ini disebut
$j$, meskipun tidak ada \myquote{j} dalam kata \myquote{imaginary}.}%
\glsadd{not:i}
\begin{equation*}
i \overset{\text{def}}{=} (0,1) .
\end{equation*}
Dengan notasi ini, $(x,y) = x+iy$.
Mulai sekarang, $x+iy$ adalah satu-satunya cara kita menulis bilangan
kompleks dalam koordinat $x$ dan $y$.
Kita menyebut $x+iy$ sebagai \emph{\myindex{bentuk Kartesius}}
bilangan kompleks.
Bilangan $i$ memiliki sifat ajaib
\begin{equation*}
i^2 = -1 .
\end{equation*}
Karena alasan ini, kadang-kadang% 
\footnote{Ada orang yang \emph{selalu} menulis $\sqrt{-1}$ alih-alih $i$.
Orang-orang itu juga pantas dipukul dengan baik.}
kita menulis $i = \sqrt{-1}$.
Perhatikan bahwa ada akar kuadrat lain dari $-1$, yaitu $-i$.
Bilangan $i$ dan $-i$ adalah solusi dari $z^2+1=0$.
Nanti akan kita buktikan bahwa setiap polinom memiliki akar dalam bilangan
kompleks.
